\subsection{学得残差的可实现补偿}
\label{sec:learned_residual_control}

前述控制比较使用完整显式模型。进一步将同一梁的$H_0$作为名义模型，把已保留的非线性磁势视作其遗漏项，使保守学习与控制补偿形成直接联系。真实仿真对象仍为$h_*$，公共参考由$H_0$在同一$u_{\rm ref}$下从零初态生成；本组参考因而不同于前述完整模型参考，误差比较限于本组对照。使用五个完整数据种子已选的GR、GC权重，估计残差
\[
 \widehat r_\theta=J\nabla H_\theta-J\nabla H_0,
 \qquad b(X)=d\cos(d^TX),
\]
以抵消所估计动量残差为目标，沿实际磁场方向求解最小二乘$\min_a\|\widehat r_{\theta,Y}+b a\|_2^2$，得到
\[
 u_\theta^{\rm ff}=-\frac{b^T\widehat r_{\theta,Y}}{b^Tb}.
\]
该负号使输入力抵消所估计的遗漏力。精确残差参照由$r_Y=-d[\sin\vartheta-\vartheta]$给出，故$u_{\rm oracle}^{\rm ff}=(\sin\vartheta-\vartheta)/\cos\vartheta$。所有方案将补偿与相同反馈相加，再对总输入施加$|u|\leq0.2$约束。无补偿组也保留同一基础反馈与参考输入。标量端口只作用于动量方向，学习误差中的运动学分量和与$b$正交的动量分量单独统计。因而状态修正$U_k\in\mathbb R^4$与标量执行器输入$u_k\in\mathbb R$有明确映射；例如步后势冲量取$U_k=(0,h b(X^f_{k+1})u_k)$，其中$X^f_{k+1}$为自由步后的配置，不将四个状态分量视为独立执行器。

原生成函数中的控制也可在这一有势实例下直接定义。对步内保持的$u_k$，在第\ref{sec:fixed_method}节中点坐标$Q=(X_{k+1}+X_k)/2$、$P=-(Y_{k+1}+Y_k)/2$中取标量截断生成量
\[
 \varphi_{h,u_k}^{[1]}(Q,P)
 =-h h_*(Q,-P)+h u_k\sin(d^TQ).
\]
令$(Y_{k+1}-Y_k,X_{k+1}-X_k)=\nabla_w\varphi_{h,u_k}^{[1]}$即得到受控中点更新。因此，该生成函数方案G与步内方案I在本限定控制势下等价；步后方案II对应实际势冲量，原对称分裂III另作辅助对照。一般向量修正要成为生成梯度还需可积条件，不能将任意$U_k$直接加到梯度上。控制从步起点采样状态后保持，固定输入的保辛结论也不直接适用于包含反馈导数的闭环映射。

在全部40条测试初值、两档步长及无补偿、精确残差、GR、GC四组中，使用相同反馈、权重与输入上限。320组独立生成函数求解与I的一步状态差不超过$5.56\times10^{-17}$，故I/G在统计中只计为同一结果。另在每条独立测试连续轨迹等距取64点，共2560个共同状态，先计算逐轨迹的分量RMS，再按每块8条轨迹及五个种子汇总。补偿后真实力残差由$3.877\times10^{-6}$降至GR的$2.927\times10^{-6}$与GC的$2.229\times10^{-6}$（表\ref{tab:control_effectiveness}左）；这些状态未参与训练或检查点选择。各自闭环经过状态上的残差另行评价，不与共同状态的估计精度混合。

共同零初态参考下，基准步长I/G的无补偿减GR杆端配对差仅为$0.02375\,\mu\mathrm m$（3/5块改善），无补偿减GC为$-0.00286\,\mu\mathrm m$（2/5改善）；精确补偿与无补偿也同处$0.7453\,\mathrm{mm}$量级。为分离初始偏差对小残差收益的影响，再让每条$H_0$名义参考与对应真实对象从相同测试初值出发，保持参考输入、反馈增益、权重及总预算不变，并对全部40条初值重复三种插入与两档步长比较。

\begin{table*}[t]
\centering\small
\caption{残差补偿的独立测试与初值匹配诊断。左：40条未训练测试轨迹各取64点的分量RMS，数值乘$10^6$；右：名义参考与真实对象从相同初值出发，I/G的杆端RMSE（$\mu\mathrm m$）。均报告五个初值块的均值$\pm$样本标准差，时间点不作为独立重复。}
\label{tab:control_effectiveness}
\begin{minipage}[t]{0.49\textwidth}\centering\setlength{\tabcolsep}{3.5pt}
(a) 共同独立测试状态\par\smallskip
\begin{tabular}{lrrr}
\toprule
补偿 & 补偿后力残差 & 运动学分量 & 正交力分量 \\
\midrule
无补偿 & $3.877\pm1.085$ & $0.000\pm0.000$ & $0.000\pm0.000$ \\
精确残差 & $0.000\pm0.000$ & $0.000\pm0.000$ & $0.000\pm0.000$ \\
GR & $2.927\pm0.465$ & $3.040\pm0.554$ & $2.069\pm0.492$ \\
GC & $2.229\pm0.683$ & $2.135\pm0.997$ & $1.458\pm0.328$ \\
\bottomrule
\end{tabular}
\end{minipage}\hfill
\begin{minipage}[t]{0.49\textwidth}\centering
(b) 初值匹配跟踪\par\smallskip
\begin{tabular}{lrr}
\toprule
补偿 & $h_0$ & $h_0/2$ \\
\midrule
无补偿 & $4.791\pm0.436$ & $1.572\pm0.137$ \\
精确残差 & $4.653\pm0.419$ & $1.379\pm0.085$ \\
GR & $4.679\pm0.427$ & $1.416\pm0.104$ \\
GC & $4.671\pm0.428$ & $1.405\pm0.101$ \\
\bottomrule
\end{tabular}
\end{minipage}
\end{table*}


初值匹配诊断中，GR和GC的杆端与完整状态误差在两档步长均较无补偿改善（5/5块）。$h_0/2$下杆端配对收益分别为$0.1568\pm0.0568$与$0.1675\pm0.0596\,\mu\mathrm m$，平均种子内相对改善为$9.85\%$与$10.52\%$。GC相对GR的杆端差仅$0.01074\pm0.00527\,\mu\mathrm m$，完整状态也只在4/5块更好。该诊断说明所选补偿能减小初值匹配时由遗漏力引发的偏差，其收益不延伸为任意参考下的稳定控制优势。

初值匹配诊断未触发饱和，峰值$|u|=0.0521$，最大补偿幅值约$2.28\times10^{-4}$；模态半径最高约$0.277$，超过学习初值上限$0.20$。旧共同参考的混合结果与这一机制诊断同时报告，结论限定于所设参考、反馈及有限时域，受控格式的数值收敛与辛结构另作检验。

